% ======================================================================
% experiments/ga.tex — 設計空間 GA: 予測 vs 実機実測
% ======================================================================

\subsection{狙い：GA の予測を実測と突き合わせる}

本レポートのここまでは\textbf{実機実測}である。一方、この一連の作業には
別の系譜がある——\texttt{aice-pi-evolution} の\textbf{設計空間 GA}
（\texttt{.aice} → Gemini レビュアー → MAP-Elites → champion）。カーネル設計を
8 軸の遺伝子（scheduler / memory\_alloc / ipc\_primitive / isolation /
arch\_target / power\_mgmt / concurrency\_safety / smp\_model）で表し、
Gemini レビュアー 4 本で採点して世代更新する。

この Round の独自性は、レビュアー rubric を\textbf{本セッションの実機実測に
基づいて}書いたことである。ゆえに GA の champion を「実測で分かっている事実」と
直接突き合わせられる——\textbf{設計空間 GA の予測 vs 実機実測}。

\subsection{結果：worker\_pool が champion（実測が裏付ける）}

seed=4 / gen=10、14 個体。MAP-Elites のセル軸は smp\_model。

\begin{table}[htbp]
\centering
\caption{GA champion（MAP-Elites セル別ベスト）}
\small
\begin{tabular}{@{}lccl@{}}
\toprule
smp\_model & best score & n & 内訳 \\
\midrule
\textbf{worker\_pool} & \textbf{0.622}（champion） & 9 & conc=lock\_free \\
full\_smp & 0.397 & 5 & conc=lock\_free \\
single\_core & (出現せず) & 0 & --- \\
\bottomrule
\end{tabular}
\end{table}

\noindent
\textbf{worker\_pool（0.622）が full\_smp（0.397）に勝った。}これは Pi 5 の
GA Round（0.507 vs 0.371）に続く再現であり、かつ Pi 3 では\textbf{実機実測が
これを裏付けた}——worker\_pool を実装して dining \spd{3.99}・nqueens
\spd{3.66} を実証している。GA の見出し（「worker\_pool を磨け、full\_smp に
走るな」）は、実測の前に、そして独立に、正しい方向を示した。

\subsection{一致とアーティファクト}

champion の 8 軸を実測と照合すると、明確な線が引ける。

\begin{table}[htbp]
\centering
\caption{GA 予測 vs 実機実測}
\small
\begin{tabular}{@{}llll@{}}
\toprule
軸 & GA champion & 実機実測 & 判定 \\
\midrule
smp\_model & worker\_pool & \spd{3.99} 実証 & $\checkmark$ 一致 \\
concurrency\_safety & lock\_free (12/14) & volatile+dsb で成立 & $\checkmark$ 一致 \\
arch\_target & arm\_rpi3 (12/14) & BCM2837 A53 & $\checkmark$ 一致 \\
isolation & mmu\_per\_process (14/14) & none/user\_kernel\_split が安全 & $\times$ アーティファクト \\
ipc\_primitive & capability/tuple\_space & sem\_mailbox（AIPL 保護） & $\times$ アーティファクト \\
\bottomrule
\end{tabular}
\end{table}

\noindent
\textbf{主軸（smp\_model・concurrency\_safety・arch）は実測と一致}する一方、
\textbf{副軸（isolation・ipc）はアーティファクト}である。全 14 個体が
\texttt{isolation=mmu\_per\_process} に固着したのはシード増殖であり、rubric が
減点しても消えなかった。\texttt{ipc=capability/tuple\_space} は実測では
AIPL/JIT/メッシュを全滅させる選択だが champion に残った。これは Pi 5 の
kernel evolution でも見た\textbf{「小規模 GA は方向性まで信頼、全遺伝子詳細は
アーティファクト混入」}パターンの再現である。

\subsection{最大の発見：支配変数が schema に無い}

本セッションの実測で最も効いた変数は\textbf{D-cache}（単一コア \spd{8.8}、
メモリ律速のスケール可否を決める、第 \ref{sec:dcache} 節）だった。しかし
8 軸の遺伝子 schema に\textbf{キャッシュ構成の軸は存在しない}。D-cache は
memory\_alloc / power\_mgmt を通じて間接的にしか表現されず、\textbf{GA は
経験的に最重要と判明した変数を構造的に探索できない}。

したがって具体的な提言が出る：schema に \texttt{cache\_policy} 軸
（\texttt{dcache\_off} / \texttt{dcache\_on\_bounce\_usb} /
\texttt{dcache\_on\_no\_dma}）を追加すべきである。そうすれば GA が D-cache の
トレードオフ（\spd{8.8} vs USB DMA の壁）を直接探索できる。

\subsection{結論：GA は方向を、実測は真実を語る}

\begin{itemize}[leftmargin=1.6em]
  \item \textbf{価値}：GA は「worker\_pool を磨け」という方向を、実測の前に
        独立に正しく示した。Pi 3 / Pi 5 で再現し、Pi 3 では実測が裏付けた。
  \item \textbf{限界 1}：全遺伝子の詳細（isolation・ipc）はアーティファクト
        混入。主軸のみ信頼でき、副軸は人手で実測に合わせて上書きすべき。
  \item \textbf{限界 2}：経験的な支配変数（D-cache）が schema に無ければ、
        GA はそれを見つけられない。実機実測が GA を補完する。
\end{itemize}
